\documentclass[conference]{IEEEtran}
\IEEEoverridecommandlockouts
\usepackage{cite}
\usepackage{amsmath,amssymb,amsfonts}
\usepackage{algorithmic}
\usepackage{graphicx}
\usepackage{textcomp}
\usepackage{xcolor}
\usepackage[utf8]{inputenc}
\usepackage[T1]{fontenc}
\usepackage[brazil]{babel}
\usepackage{booktabs}
\usepackage{array}

\def\BibTeX{{\rm B\kern-.05em{\sc i\kern-.025em b}\kern-.08em
    T\kern-.1667em\lower.7ex\hbox{E}\kern-.125emX}}

\begin{document}

\title{Player de Música com Requisitos de Tempo Real\\
{\footnotesize Laboratório 01 — Disciplina: Processamento em Tempo Real — 1º/2026}
}

\author{
\IEEEauthorblockN{Monique Yafuso Piquet}
\IEEEauthorblockA{
    Matrícula: 17/0172091 \\
    \textit{Engenharia Mecatrônica - Controle e Automação} \\
    \textit{Universidade de Brasília} \\
    Brasília, Brasil \\
    piquetmy@gmail.com
}
\and
\IEEEauthorblockN{Raphael Henrique Nogueira Silva}
\IEEEauthorblockA{
    Matrícula: 22/2005930 \\
    \textit{Engenharia Mecatrônica - Controle e Automação} \\
    \textit{Universidade de Brasília} \\
    Brasília, Brasil \\
    raphaelhenriqueas@gmail.com
}
\and
\IEEEauthorblockN{Thiago Carvalho Silva}
\IEEEauthorblockA{
    Matrícula: 22/2006463 \\
    \textit{Engenharia Mecatrônica - Controle e Automação} \\
    \textit{Universidade de Brasília} \\
    Brasília, Brasil \\
    thiago.carvalho861@hotmail.com
}
}

\maketitle

% ============================================================
\begin{abstract}
Este trabalho apresenta o desenvolvimento de um tocador de música
portátil inspirado no iPod clássico, implementado sobre a plataforma
ESP32. O sistema integra um display LCD 16x2 para interface com o
usuário, um módulo leitor de cartão SD para armazenamento de arquivos
de áudio, botões de controle, codificador rotativo para navegação e volume,
e um conversor digital/analógico (DAC) externo via protocolo I2S para
saída de áudio de alta qualidade. O desenvolvimento é realizado em lin-
guagem Python, com ênfase nos conceitos de sistemas de tempo real, incluindo
gerenciamento de tarefas, controle de buffers e tratamento de interrupções.
Espera-se que o sistema realize a reprodução contínua e sem falhas de
arquivos de áudio armazenados no cartão SD, respondendo de forma de-
terminística aos comandos do usuário.
\end{abstract}

% ============================================================
\section{Introdução}
Um Sistema de Tempo Real é definido como aquele em que a correção do sistema depende não
apenas do resultado lógico, mas também do instante em que
esses resultados são produzidos. Nesse contexto, o tocador de
música representa uma aplicação com restrições temporais bem definidas: a
tarefa de reprodução de áudio deve entregar amostras ao DAC a uma taxa
constante de, tipicamente, 44.1 kHz, sob pena de introduzir artefatos
sonoros audíveis, caracterizando uma violação de deadline.

Este projeto propõe o desenvolvimento de um tocador de música portátil baseado
no microcontrolador ESP32. O sistema
é composto por um display LCD 16x2 para exibição de metadados e estado de reprodução; leitor
de cartão SD para acesso a arquivos MP3, além de arquivos que facilitem a indexação; botões físicos de controle (play/pause,
próxima faixa, faixa anterior); codificador rotativo para ajuste de volume e
navegação; DAC externo 16 bits via I2S para saída de áudio. O sistema é construído
à partir do microkernel FreeRTOS, utilizando tarefas periódicas e aperiódicas e interrupções de hardware, atendendo aos
requisitos formais de um STR.

% ============================================================
\section{Objetivos}

O objetivo geral deste trabalho é desenvolver e validar um tocador de música
embarcado que respeite os requisitos temporais e funcionais de um sistema de
tempo real. Os objetivos específicos são:

\begin{itemize}
    \item Definir e especificar o modelo de tarefas do sistema, incluindo
    períodos e deadlines;
    \item Implementar tarefas periódicas e aperiódicas/esporádicas com
    tratamento de interrupções de software e hardware;
    \item Validar o sistema por meio de simulação no ambiente Wokwi, como posterior implementação física;
    \item Demonstrar que a aplicação satisfaz os deadlines estabelecidos sem
    violações perceptíveis.
\end{itemize}

% ============================================================
\section{Fundamentação Teórica}

\subsection{Sistemas de Tempo Real}

Um sistema de tempo real é aquele cujo comportamento correto depende tanto
da lógica das operações quanto dos instantes em que seus resultados são
produzidos. As tarefas em um STR são classificadas segundo
seus deadlines em: (i) \textit{hard real-time}, em que a perda de um deadline
representa falha catastrófica; (ii) \textit{firm real-time}, em que resultados
tardios são descartados sem valor; e (iii) \textit{soft real-time}, em que
atrasos degradam a qualidade mas não comprometem a operação.

No tocador de música, a tarefa de envio de amostras ao DAC é classificada como
\textit{firm real-time}: uma amostra entregue fora do instante correto é
simplesmente descartada, gerando artefato sonoro. As demais tarefas (leitura SD,
atualização de display, tratamento de botões...) são \textit{soft real-time}.

\subsection{Modelo de Tarefas}

Cada tarefa $\tau_i$ é definida por
$C_i, T_i, D_i, J_i$, onde $C_i$ é o tempo de computação no pior caso
(WCET), $T_i$ é o período, $D_i$ é o deadline relativo e $J_i$ é o jitter máximo tolerado.

\subsection{FreeRTOS e ESP-ADF}

O FreeRTOS é um microkernel de tempo real, amplamente
utilizado em sistemas embarcados. Sobre o ESP32, ele é
integrado ao ESP-IDF (Espressif IoT Development Framework) e oferece:
escalonador preemptivo por prioridades, semáforos binários e de contagem,
mutexes, filas (queues) e notificações de tarefas. O ESP32 possui dois núcleos
de 240 MHz, o que permite fixar a tarefa crítica de áudio em um núcleo
dedicado (Core 1) e as demais tarefas no Core 0, eliminando interferências. Ainda, a Espressif dispõe do framework de áudio ESP-ADF, pelo qual é possível a decodificação de arquivos em MP3, bem como o controle de áudio em buffer, para enfim seu \textit{output} via protocolo I2S.

\subsection{Protocolo I2S e DAC Externo}

O I2S (\textit{Inter-IC Sound}) é um barramento serial síncrono desenvolvido
para interconexão de circuitos de áudio digital. Utiliza três
sinais: SCK, WS e
SD (\textit{serial data}). O DAC externo recebe as amostras decodificadas do ESP32 e as converte
em tensão analógica com maior resolução e melhor resolução do que o DAC interno do ESP32.

\subsection{Buffer Circular}

O buffer circular é a estrutura de dados central que desacopla a tarefa de
leitura do SD (produtora) da tarefa de reprodução de áudio (consumidora). Seu
tamanho mínimo é dado por:

\begin{equation}
    B_{min} = T_{lat} \times F_s \times N_b
    \label{eq_buffer}
\end{equation}

onde $T_{lat}$ é a latência máxima esperada na leitura do SD (s), $F_s$ é a
frequência de amostragem (Hz) e $N_b$ é o número de bytes por amostra. Para
$T_{lat} = 10\,\text{ms}$, $F_s = 44100\,\text{Hz}$ e $N_b = 4$
(estéreo, 16 bits), obtém-se $B_{min} \approx 1764\,\text{bytes}$. Na
implementação, adota-se $B = 4096\,\text{bytes}$ para margem de segurança.

% ============================================================
\section{Materiais e Métodos}

\subsection{Plataforma de Hardware}

\subsubsection{ESP32}
O ESP32 (Espressif Systems) é um SoC de duplo núcleo Xtensa LX6 com clock de 240 MHz, 520 KB de SRAM, suporte nativo a SPI, I2C e I2S,
e integração com o FreeRTOS via ESP-IDF.

\subsubsection{Display LCD 16x2 com I2C}
O display LCD HD44780 permite comunicação via 8 bits em paralelo (mais eficiente que I2C nesse caso), exibindo nome da faixa, estado
do player e nível de volume.

\subsubsection{Módulo Leitor de Cartão SD (SPI)}
O módulo cartão SD comunica-se via SPI (MOSI, MISO, SCK, CS). O
sistema de arquivos utilizado é FAT32, acessado pela biblioteca
\texttt{esp\_vfs\_fat} do ESP-IDF. Os arquivos de áudio são armazenados no
formato MP3, 44100 Hz, estéreo, 256kbps.

\subsubsection{Botões de Controle}
Três botões físicos são conectados a pinos GPIO com resistores de 1k$\omega$
externos: LEFT (diminuir volume), MID (play/pause, anterior, próxima, stop) e RIGHT
(aumentar volume). Cada botão é configurado para gerar interrupção na borda
de descida (\texttt{GPIO\_INTR\_NEGEDGE}).

\subsubsection{Codificador Rotativo}
O codificador rotativo fornece dois sinais em quadratura (CLK e DT)
e um sinal de botão (SW). Os pinos CLK e DT são configurados com interrupção
em ambas as bordas para detecção de direção de giro. O encoder é utilizado
para navegação (giro) e confirmação de seleção (botão).

\subsubsection{DAC Externo via I2S}
Um DAC externo compatível com I2S (ex.: PCM5102A) é conectado ao ESP32 pelos
pinos I2S\_BCK, I2S\_WS e I2S\_DATA. O driver I2S do ESP-IDF é configurado
no modo mestre com taxa de amostragem de 44.100 Hz, formato 16 bits estéreo.
